\documentclass[11pt,a4paper]{article}

\usepackage{iftex}
\ifPDFTeX
  \usepackage[T1]{fontenc}
  \usepackage[utf8]{inputenc}
  \usepackage{lmodern}
\else
  \usepackage{fontspec}
\fi

\usepackage[english]{babel}
\usepackage[margin=2.5cm]{geometry}
\usepackage{amsmath,amssymb}
\usepackage{array}
\usepackage{booktabs}
\usepackage{enumitem}
\usepackage{hyperref}

\hypersetup{
  colorlinks=true,
  linkcolor=black,
  urlcolor=blue,
  pdftitle={Genetic algorithm for classification},
  pdfauthor={Fellipe Almassy}
}

\newcommand{\cod}[1]{\texttt{#1}}
\newcommand{\chrom}{\mathbf{c}}

\title{Genetic algorithm for classification}
\author{Fellipe Almassy}
\date{November 2026}

\begin{document}

\maketitle


\begin{abstract}
This document explains how a genetic algorithm can train a linear classifier.
Instead of using gradients, the method evolves a population of candidate solutions:
the best ones are more likely to reproduce, their offspring are mutated and replace
the worst solutions. The text describes each step of the algorithm, with the
mathematical notation and a numerical example, and points to the matching function
of the \cod{genetic-classifier} library.
\end{abstract}

\section{Goal}

We start from a labeled dataset made of samples. Each sample has a vector of
\emph{features} and a known class. The goal of the genetic algorithm is to find the
coefficients of a linear function that, from the features, tells the class of each
sample. After training, this function can classify new data.

\section{Notation}

\begin{center}
\begin{tabular}{cl}
\toprule
Symbol & Meaning \\
\midrule
$N$ & number of samples in the dataset \\
$n$ & number of features \\
$K$ & number of classes, numbered from $0$ to $K-1$ \\
$\mathbf{x} = (x_1, \dots, x_n)$ & features of a sample \\
$y$ & true class of a sample \\
$m$ & number of chromosomes in the population \\
$L = n + 1$ & number of genes in each chromosome \\
$h$ & number of offspring created per generation \\
$r$ & number of genes changed in each child \\
$s$ & number of chromosomes replaced per generation \\
\bottomrule
\end{tabular}
\end{center}

\section{Representing a solution}

Each \textbf{chromosome} is a vector of real numbers that represents a candidate
solution. It is made of a \textbf{bias} $b$ and one \textbf{gene} $g_j$ for each
feature:
\begin{equation}
  \chrom = (b, g_1, g_2, \dots, g_n).
\end{equation}

With these values, the chromosome computes a score $q$ for any sample:
\begin{equation}
  q(\mathbf{x}) = b + g_1 x_1 + g_2 x_2 + \dots + g_n x_n
                = b + \sum_{j=1}^{n} g_j x_j.
  \label{eq:score}
\end{equation}

Each gene is the weight of a feature in the decision. The bias shifts the score of
every sample by the same amount. Without it, a sample with all features equal to zero
would always get $q = 0$, and the decision boundary would be forced to pass through
the origin.

In binary classification, the score is turned into a class by its sign:
\begin{equation}
  \hat{y} =
  \begin{cases}
    1, & \text{if } q(\mathbf{x}) \geq 0,\\
    0, & \text{if } q(\mathbf{x}) < 0.
  \end{cases}
  \label{eq:binary}
\end{equation}

In ordinal classification, the $K$ classes have an order (for example,
bad $<$ average $<$ good). We define $K-1$ fixed cut points
$t_1 < t_2 < \dots < t_{K-1}$, chosen by the user, which do not evolve during
training. With the conventions $t_0 = -\infty$ and $t_K = +\infty$, the predicted
class is
\begin{equation}
  \hat{y} = k \quad \text{such that} \quad t_k \leq q(\mathbf{x}) < t_{k+1}.
  \label{eq:ordinal}
\end{equation}

\section{The algorithm loop}

\begin{enumerate}[nosep]
  \item Create the initial population.
  \item Compute the fitness of each chromosome.
  \item Select a father and a mother by roulette.
  \item Create $h$ offspring by crossover.
  \item Mutate the offspring.
  \item Replace the $s$ worst chromosomes with the $s$ best offspring and go back to step~2.
\end{enumerate}

\medskip
The loop ends when the best fitness found reaches the target value or when the
maximum number of generations runs out. The result is the best chromosome seen during
the whole training, kept aside, and not necessarily a chromosome of the last
population.

\subsection*{Step 1: creating the population}

The initial population is a matrix $P$ with $m$ rows and $L$ columns. Each row is a
chromosome and each column is a gene, the first column being the bias:
\begin{equation}
  P =
  \begin{pmatrix}
    b_1    & g_{1,1} & g_{1,2} & \cdots & g_{1,n} \\
    b_2    & g_{2,1} & g_{2,2} & \cdots & g_{2,n} \\
    \vdots & \vdots  & \vdots  & \ddots & \vdots  \\
    b_m    & g_{m,1} & g_{m,2} & \cdots & g_{m,n}
  \end{pmatrix}.
\end{equation}

All values are drawn uniformly from an interval, $[-1, 1]$ by default. The limits are
parameters of the function and can be changed.

\subsection*{Step 2: computing the fitness}

The fitness is the grade of a chromosome: the better it classifies the samples of the
dataset, the higher the grade. The most direct measure would be accuracy, that is,
the fraction of correct answers. It is misleading, however, when the classes are
imbalanced. In a dataset with 95 samples of class 1 and 5 of class 0, a chromosome
that always answers 1 has 95\% accuracy without having learned anything.

That is why we measure the recall of each class, that is, the fraction of samples of
that class that the chromosome got right:
\begin{equation}
  R_k = \frac{\text{correct answers in class } k}{\text{total samples of class } k}.
\end{equation}

The fitness is the product of the recalls:
\begin{equation}
  F = \prod_{k=0}^{K-1} R_k.
  \label{eq:fitness}
\end{equation}

So $0 \leq F \leq 1$, and $F = 0$ whenever the chromosome gets every sample of some
class wrong. In the case above, the chromosome that always answers 1 has $R_1 = 1$ and
$R_0 = 0$, so $F = 0$. The result depends only on the fraction of correct answers
within each class, not on how many samples each class has.

\paragraph{Example.} Consider the dataset in Table~\ref{tab:data} and the chromosome
$\chrom = (2,\, 3,\, -4,\, 1,\, -1)$. The values are integers and outside the default
interval only to make the arithmetic easier. The score is
\begin{equation*}
  q(\mathbf{x}) = 2 + 3x_1 - 4x_2 + 1x_3 - 1x_4.
\end{equation*}

\begin{table}[ht]
\centering
\caption{Example dataset.}
\label{tab:data}
\begin{tabular}{lccccc}
\toprule
 & $x_1$ & $x_2$ & $x_3$ & $x_4$ & Class \\
\midrule
Sample 1 & 1 & 0 & 1 & 1 & 1 \\
Sample 2 & 1 & 1 & 0 & 0 & 1 \\
Sample 3 & 1 & 0 & 0 & 1 & 0 \\
Sample 4 & 0 & 1 & 1 & 1 & 0 \\
\bottomrule
\end{tabular}
\end{table}

Table~\ref{tab:evaluation} applies Equation~\eqref{eq:score} to each sample and turns
the score into a class with Equation~\eqref{eq:binary}.

\begin{table}[ht]
\centering
\caption{Evaluating the chromosome on each sample.}
\label{tab:evaluation}
\begin{tabular}{llcccc}
\toprule
 & Computation & $q$ & Predicted & True & Result \\
\midrule
Sample 1 & $2 + 3\cdot1 - 4\cdot0 + 1\cdot1 - 1\cdot1$ & $5$  & 1 & 1 & correct \\
Sample 2 & $2 + 3\cdot1 - 4\cdot1 + 1\cdot0 - 1\cdot0$ & $1$  & 1 & 1 & correct \\
Sample 3 & $2 + 3\cdot1 - 4\cdot0 + 1\cdot0 - 1\cdot1$ & $4$  & 1 & 0 & wrong \\
Sample 4 & $2 + 3\cdot0 - 4\cdot1 + 1\cdot1 - 1\cdot1$ & $-2$ & 0 & 0 & correct \\
\bottomrule
\end{tabular}
\end{table}

The chromosome got 3 of the 4 samples right, an accuracy of 75\%. Per class,
\begin{equation*}
  R_1 = \frac{2}{2} = 1, \qquad R_0 = \frac{1}{2} = 0.5, \qquad
  F = R_1 \cdot R_0 = 0.5.
\end{equation*}

\subsection*{Step 3: roulette selection}

The fitnesses are normalized into probabilities that add up to 1:
\begin{equation}
  p_i = \frac{F_i}{\sum_{l=1}^{m} F_l}.
\end{equation}
If every fitness is zero, every chromosome gets the same chance, $p_i = 1/m$.

To spin the roulette, we draw $u$ uniformly in $[0, 1)$ and pick the first chromosome
whose cumulative probability exceeds $u$. The draw is done twice, once for the father
and once for the mother, so both can be the same chromosome. A single couple is chosen
per generation.

\paragraph{Example.} With $F = (0.5,\ 0.25,\ 0.25)$, we get
$p = (0.5,\ 0.25,\ 0.25)$, and the cumulative probabilities are
$(0.5,\ 0.75,\ 1)$. If $u = 0.6$, the second chromosome is picked.

\subsection*{Step 4: one-point crossover}

For each of the $h$ offspring, we draw a cut point $c \in \{1, \dots, L-1\}$. The
child gets the first $c$ genes from the father and the remaining genes from the
mother:
\begin{equation}
  \chrom^{\text{child}} =
  (\chrom^{\text{father}}_{1}, \dots, \chrom^{\text{father}}_{c},\;
   \chrom^{\text{mother}}_{c+1}, \dots, \chrom^{\text{mother}}_{L}).
\end{equation}
Since the cut is never at the ends, the child always inherits at least one gene from
each parent.

\paragraph{Example.} With father $(2,\, 3,\, -4,\, 1,\, -1)$, mother
$(-1,\, 0.5,\, 2,\, -3,\, 1)$ and $c = 2$, the child is $(2,\, 3,\, 2,\, -3,\, 1)$.

\subsection*{Step 5: mutation}

In each child, we draw $r$ different positions and replace the gene at each one with
a new uniform value in $[-1, 1]$. Mutation brings into the population values that did
not exist in the parents and therefore helps keep diversity. In the current
implementation, this interval is fixed and does not follow the limits chosen in
step~1.

\paragraph{Example.} If the drawn position is the fourth gene and the new value is
$0.4$, the child $(2,\, 3,\, 2,\, -3,\, 1)$ becomes $(2,\, 3,\, 2,\, 0.4,\, 1)$.

\subsection*{Step 6: updating the population}

We compute the fitness of the mutated offspring. The $s$ chromosomes with the lowest
fitness leave the population and the $s$ offspring with the highest fitness take
their place, with $0 \leq s \leq \min(m, h)$. The rest of the population moves
unchanged to the next generation, and the remaining offspring are discarded.

The swap does not compare children and parents: a child can enter even with a lower
fitness than the chromosome it replaced. When $s < m$, the best chromosomes of the
population are preserved, since only the worst ones leave.

\section{Parameters}

Table~\ref{tab:parameters} maps the symbols of this text to the parameters of the
\cod{train} function.

\begin{table}[ht]
\centering
\caption{Parameters of \cod{train} and their default values.}
\label{tab:parameters}
\begin{tabular}{clc}
\toprule
Symbol & Parameter & Default \\
\midrule
$m$ & \cod{n\_chromosomes} & 30 \\
$r$ & \cod{n\_mutated\_genes} & 2 \\
$h$ & \cod{n\_offspring} & 16 \\
$s$ & \cod{n\_replacements} & half of $h$ (8) \\
--- & \cod{n\_generations} & 1000 \\
--- & \cod{target\_fitness} & 1.0 \\
$t_1, \dots, t_{K-1}$ & \cod{thresholds} & none (binary) \\
--- & \cod{gene\_min}, \cod{gene\_max} & $-1$ and $1$ \\
\bottomrule
\end{tabular}
\end{table}

\section{Where each step is in the code}

\begin{table}[ht]
\centering
\caption{Where each step is implemented, in \cod{src/genetic\_classifier/}.}
\begin{tabular}{lll}
\toprule
Step & File & Function \\
\midrule
1. Creation    & \cod{creation.py}   & \cod{create\_chromosomes} \\
2. Fitness     & \cod{fitness.py}    & \cod{binary\_fitness}, \cod{ordinal\_fitness} \\
3. Selection   & \cod{fitness.py}    & \cod{fitness\_proportions} \\
               & \cod{selection.py}  & \cod{select\_parents\_roulette} \\
4. Crossover   & \cod{crossover.py}  & \cod{one\_point\_crossover} \\
5. Mutation    & \cod{mutation.py}   & \cod{mutate} \\
6. Update      & \cod{update.py}     & \cod{update\_population} \\
Full loop      & \cod{core.py}       & \cod{train} \\
Classification & \cod{core.py}       & \cod{predict} \\
\bottomrule
\end{tabular}
\end{table}

\end{document}
